%% ma: L10-B14-0066
%% lop: 10
%% chuong: 5
%% bai: 14
%% dang: 10.14.9
%% loai: TN4
%% muc: VD
%% dap_an: "B"
%% nguon: "chuyen-de-cac-so-dac-trung-cua-mau-so-lieu-khong-ghep-nhom-toan-10-chuong-trinh-moi.pdf (GV. Phan Nhật Linh) (P05), Bài 03 Dạng 2 Phần I Câu 10"
%% kien_thuc: "Lớp 10 Bài 13 (số trung bình), Bài 14 (phương sai, độ lệch chuẩn)"
%% trang_thai: chinh-thuc
%% ghi_chu: "Tự duyệt theo yêu cầu giáo viên 22/10/2026."
\begin{ex}
Điểm trung bình các môn học của hai bạn An và Bình trong năm học qua được cho trong bảng sau:
\begin{center}
\begin{tabular}{|c|c|c|c|c|c|c|c|c|}\hline
 & Toán & Lí & Hóa & Sinh & Văn & Sử & Địa & Tiếng Anh\\\hline
An & $8{,}0$ & $6{,}5$ & $8{,}8$ & $8{,}3$ & $8{,}0$ & $8{,}0$ & $7{,}2$ & $9{,}0$\\\hline
Bình & $8{,}5$ & $9{,}0$ & $7{,}5$ & $8{,}0$ & $6{,}5$ & $8{,}5$ & $7{,}5$ & $7{,}0$\\\hline
\end{tabular}
\end{center}
Căn cứ vào phương sai, bạn nào học ``lệch'' hơn?
\choice{An}{\True Bình}{Mức độ học ``lệch'' của hai bạn như nhau}{Không đủ cơ sở để kết luận}
\loigiai{An: $\overline x=7{,}975$; $s^2=\dfrac18\sum x_i^2-\overline x^2=\dfrac{923}{1600}\approx0{,}58$.

Bình: $\overline x=7{,}8125$; $s^2=\dfrac{159}{256}\approx0{,}62$.

Phương sai điểm của Bình lớn hơn nên Bình học lệch hơn. Chọn B.}
\end{ex}
